% Auto-generated mixed-window report (Arabidopsis/Rice 0 bp; others 25 kb)

\subsection{Cross-species genome-wide associations with soil nitrogen (0--5 cm) identify shared orthologous genes across all species}

We performed GWAS in Arabidopsis, barley, rice, maize, and sorghum using a mixed linear model and kinship correction in GEMMA. Under Bonferroni correction, the number of significant SNPs was heterogeneous across species (Arabidopsis 64, rice 202, sorghum 2, barley 4, maize 1). To evaluate cross-species convergence beyond the Bonferroni tail, we selected the top 0.5\% of SNPs per species and mapped nearby genes to orthogroups, using gene-body-only annotation (0 bp) for Arabidopsis and rice and \pm25 kb annotation for maize, sorghum, and barley. This yielded 2,690 Arabidopsis genes, 2,844 rice genes, 4,423 maize genes, 1,636 barley genes, and 990 sorghum genes.

Orthogroup integration identified 21,316 total orthogroups, with 0 orthogroups represented in all five species. Table~\ref{tab:shared_function_candidates_soilN_mixed} summarizes these shared orthogroups and reports, for each species, the strongest gene-level signal within each orthogroup as $-\log_{10}(P)$.

\begin{table*}[t]
\caption{Shared orthogroups across all five species from the top 0.5\% of SNPs for soil nitrogen (0--5 cm) (mixed-window annotation).\label{tab:shared_function_candidates_soilN_mixed}}
\scriptsize
\tabcolsep=3pt
\begin{tabular*}{\textwidth}{@{\extracolsep{\fill}}lllllllllll@{}}
\toprule
Orthogroup & Maize gene & $-\log_{10}(P)$ & Sorghum gene & $-\log_{10}(P)$ & Rice gene & $-\log_{10}(P)$ & Barley gene & $-\log_{10}(P)$ & Arabidopsis gene & $-\log_{10}(P)$ \\
\midrule
NA & NA & NA & NA & NA & NA & NA & NA & NA & NA & NA \\
\botrule
\end{tabular*}
\begin{tablenotes}%
\item Note: One representative gene per species is shown for each orthogroup (the gene with the smallest $P$ within that species/orthogroup).
\end{tablenotes}
\end{table*}

\subsection{Cross-species genome-wide associations with WorldClim PC1 identify shared orthologous genes across all species}

We performed GWAS in Arabidopsis, barley, rice, maize, and sorghum using a mixed linear model and kinship correction in GEMMA. Under Bonferroni correction, the number of significant SNPs was heterogeneous across species (Arabidopsis 38, rice 602, sorghum 6, barley 8, maize 1). To evaluate cross-species convergence beyond the Bonferroni tail, we selected the top 0.5\% of SNPs per species and mapped nearby genes to orthogroups, using gene-body-only annotation (0 bp) for Arabidopsis and rice and \pm25 kb annotation for maize, sorghum, and barley. This yielded 3,185 Arabidopsis genes, 4,640 rice genes, 3,852 maize genes, 1,270 barley genes, and 1,283 sorghum genes.

Orthogroup integration identified 21,013 total orthogroups, with 0 orthogroups represented in all five species. Table~\ref{tab:shared_function_candidates_pc1_worldclim_mixed} summarizes these shared orthogroups and reports, for each species, the strongest gene-level signal within each orthogroup as $-\log_{10}(P)$.

\begin{table*}[t]
\caption{Shared orthogroups across all five species from the top 0.5\% of SNPs for WorldClim PC1 (mixed-window annotation).\label{tab:shared_function_candidates_pc1_worldclim_mixed}}
\scriptsize
\tabcolsep=3pt
\begin{tabular*}{\textwidth}{@{\extracolsep{\fill}}lllllllllll@{}}
\toprule
Orthogroup & Maize gene & $-\log_{10}(P)$ & Sorghum gene & $-\log_{10}(P)$ & Rice gene & $-\log_{10}(P)$ & Barley gene & $-\log_{10}(P)$ & Arabidopsis gene & $-\log_{10}(P)$ \\
\midrule
NA & NA & NA & NA & NA & NA & NA & NA & NA & NA & NA \\
\botrule
\end{tabular*}
\begin{tablenotes}%
\item Note: One representative gene per species is shown for each orthogroup (the gene with the smallest $P$ within that species/orthogroup).
\end{tablenotes}
\end{table*}

\subsection{Cross-species genome-wide associations with WorldClim PC2 identify shared orthologous genes across all species}

We performed GWAS in Arabidopsis, barley, rice, maize, and sorghum using a mixed linear model and kinship correction in GEMMA. Under Bonferroni correction, the number of significant SNPs was heterogeneous across species (Arabidopsis 13, rice 0, sorghum 0, barley 1, maize 1). To evaluate cross-species convergence beyond the Bonferroni tail, we selected the top 0.5\% of SNPs per species and mapped nearby genes to orthogroups, using gene-body-only annotation (0 bp) for Arabidopsis and rice and \pm25 kb annotation for maize, sorghum, and barley. This yielded 2,752 Arabidopsis genes, 3,745 rice genes, 3,953 maize genes, 1,557 barley genes, and 1,151 sorghum genes.

Orthogroup integration identified 20,931 total orthogroups, with 0 orthogroups represented in all five species. Table~\ref{tab:shared_function_candidates_pc2_worldclim_mixed} summarizes these shared orthogroups and reports, for each species, the strongest gene-level signal within each orthogroup as $-\log_{10}(P)$.

\begin{table*}[t]
\caption{Shared orthogroups across all five species from the top 0.5\% of SNPs for WorldClim PC2 (mixed-window annotation).\label{tab:shared_function_candidates_pc2_worldclim_mixed}}
\scriptsize
\tabcolsep=3pt
\begin{tabular*}{\textwidth}{@{\extracolsep{\fill}}lllllllllll@{}}
\toprule
Orthogroup & Maize gene & $-\log_{10}(P)$ & Sorghum gene & $-\log_{10}(P)$ & Rice gene & $-\log_{10}(P)$ & Barley gene & $-\log_{10}(P)$ & Arabidopsis gene & $-\log_{10}(P)$ \\
\midrule
NA & NA & NA & NA & NA & NA & NA & NA & NA & NA & NA \\
\botrule
\end{tabular*}
\begin{tablenotes}%
\item Note: One representative gene per species is shown for each orthogroup (the gene with the smallest $P$ within that species/orthogroup).
\end{tablenotes}
\end{table*}

\subsection{Cross-species genome-wide associations with WorldClim PC3 identify shared orthologous genes across all species}

We performed GWAS in Arabidopsis, barley, rice, maize, and sorghum using a mixed linear model and kinship correction in GEMMA. Under Bonferroni correction, the number of significant SNPs was heterogeneous across species (Arabidopsis 29, rice 10, sorghum 80, barley 24, maize 0). To evaluate cross-species convergence beyond the Bonferroni tail, we selected the top 0.5\% of SNPs per species and mapped nearby genes to orthogroups, using gene-body-only annotation (0 bp) for Arabidopsis and rice and \pm25 kb annotation for maize, sorghum, and barley. This yielded 2,904 Arabidopsis genes, 4,363 rice genes, 3,961 maize genes, 1,629 barley genes, and 554 sorghum genes.

Orthogroup integration identified 20,201 total orthogroups, with 0 orthogroups represented in all five species. Table~\ref{tab:shared_function_candidates_pc3_worldclim_mixed} summarizes these shared orthogroups and reports, for each species, the strongest gene-level signal within each orthogroup as $-\log_{10}(P)$.

\begin{table*}[t]
\caption{Shared orthogroups across all five species from the top 0.5\% of SNPs for WorldClim PC3 (mixed-window annotation).\label{tab:shared_function_candidates_pc3_worldclim_mixed}}
\scriptsize
\tabcolsep=3pt
\begin{tabular*}{\textwidth}{@{\extracolsep{\fill}}lllllllllll@{}}
\toprule
Orthogroup & Maize gene & $-\log_{10}(P)$ & Sorghum gene & $-\log_{10}(P)$ & Rice gene & $-\log_{10}(P)$ & Barley gene & $-\log_{10}(P)$ & Arabidopsis gene & $-\log_{10}(P)$ \\
\midrule
NA & NA & NA & NA & NA & NA & NA & NA & NA & NA & NA \\
\botrule
\end{tabular*}
\begin{tablenotes}%
\item Note: One representative gene per species is shown for each orthogroup (the gene with the smallest $P$ within that species/orthogroup).
\end{tablenotes}
\end{table*}

\subsection{Cross-species genome-wide associations with soil pH identify shared orthologous genes across all species}

We performed GWAS in Arabidopsis, barley, rice, maize, and sorghum using a mixed linear model and kinship correction in GEMMA. Under Bonferroni correction, the number of significant SNPs was heterogeneous across species (Arabidopsis 6, rice 383, sorghum 3, barley 6, maize 0). To evaluate cross-species convergence beyond the Bonferroni tail, we selected the top 0.5\% of SNPs per species and mapped nearby genes to orthogroups, using gene-body-only annotation (0 bp) for Arabidopsis and rice and \pm25 kb annotation for maize, sorghum, and barley. This yielded 2,800 Arabidopsis genes, 5,104 rice genes, 3,937 maize genes, 1,526 barley genes, and 1,203 sorghum genes.

Orthogroup integration identified 21,152 total orthogroups, with 4 orthogroups represented in all five species. Table~\ref{tab:shared_function_candidates_ph_mixed} summarizes these shared orthogroups and reports, for each species, the strongest gene-level signal within each orthogroup as $-\log_{10}(P)$.

\begin{table*}[t]
\caption{Shared orthogroups across all five species from the top 0.5\% of SNPs for soil pH (mixed-window annotation).\label{tab:shared_function_candidates_ph_mixed}}
\scriptsize
\tabcolsep=3pt
\begin{tabular*}{\textwidth}{@{\extracolsep{\fill}}lllllllllll@{}}
\toprule
Orthogroup & Maize gene & $-\log_{10}(P)$ & Sorghum gene & $-\log_{10}(P)$ & Rice gene & $-\log_{10}(P)$ & Barley gene & $-\log_{10}(P)$ & Arabidopsis gene & $-\log_{10}(P)$ \\
\midrule
179088at33090 & Zm00001eb254960 & 3.19 & SORBI\_3003G341500 & 2.54 & LOC\_Os07g01010 & 2.89 & HORVU.MOREX.r3.5HG0465270 & 2.82 & AT5G60020 & 2.86 \\
220610at3193 & Zm00001eb254960 & 3.19 & SORBI\_3003G341500 & 2.54 & LOC\_Os07g01010 & 2.89 & HORVU.MOREX.r3.5HG0465270 & 2.82 & AT5G60020 & 2.86 \\
4495752at2759 & Zm00001eb254960 & 3.19 & SORBI\_3003G341500 & 2.54 & LOC\_Os07g01010 & 2.89 & HORVU.MOREX.r3.5HG0465270 & 2.82 & AT5G60020 & 2.86 \\
640662at2759 & Zm00001eb068210 & 3.57 & SORBI\_3003G320800 & 3.31 & LOC\_Os05g01620 & 3.00 & HORVU.MOREX.r3.2HG0112720 & 3.53 & AT3G17070 & 2.62 \\
\botrule
\end{tabular*}
\begin{tablenotes}%
\item Note: One representative gene per species is shown for each orthogroup (the gene with the smallest $P$ within that species/orthogroup).
\end{tablenotes}
\end{table*}

\subsection{Cross-species genome-wide associations with AM fungal relative abundance colonization identify shared orthologous genes across all species}

We performed GWAS in Arabidopsis, barley, rice, maize, and sorghum using a mixed linear model and kinship correction in GEMMA. Under Bonferroni correction, the number of significant SNPs was heterogeneous across species (Arabidopsis 824, rice 268, sorghum 0, barley 9, maize 1). To evaluate cross-species convergence beyond the Bonferroni tail, we selected the top 0.5\% of SNPs per species and mapped nearby genes to orthogroups, using gene-body-only annotation (0 bp) for Arabidopsis and rice and \pm25 kb annotation for maize, sorghum, and barley. This yielded 3,385 Arabidopsis genes, 3,827 rice genes, 3,791 maize genes, 1,246 barley genes, and 994 sorghum genes.

Orthogroup integration identified 20,577 total orthogroups, with 0 orthogroups represented in all five species. Table~\ref{tab:shared_function_candidates_am_rel_abundance_colonization_mixed} summarizes these shared orthogroups and reports, for each species, the strongest gene-level signal within each orthogroup as $-\log_{10}(P)$.

\begin{table*}[t]
\caption{Shared orthogroups across all five species from the top 0.5\% of SNPs for AM fungal relative abundance colonization (mixed-window annotation).\label{tab:shared_function_candidates_am_rel_abundance_colonization_mixed}}
\scriptsize
\tabcolsep=3pt
\begin{tabular*}{\textwidth}{@{\extracolsep{\fill}}lllllllllll@{}}
\toprule
Orthogroup & Maize gene & $-\log_{10}(P)$ & Sorghum gene & $-\log_{10}(P)$ & Rice gene & $-\log_{10}(P)$ & Barley gene & $-\log_{10}(P)$ & Arabidopsis gene & $-\log_{10}(P)$ \\
\midrule
NA & NA & NA & NA & NA & NA & NA & NA & NA & NA & NA \\
\botrule
\end{tabular*}
\begin{tablenotes}%
\item Note: One representative gene per species is shown for each orthogroup (the gene with the smallest $P$ within that species/orthogroup).
\end{tablenotes}
\end{table*}

\subsection{Cross-species genome-wide associations with AM fungal roots colonized identify shared orthologous genes across all species}

We performed GWAS in Arabidopsis, barley, rice, maize, and sorghum using a mixed linear model and kinship correction in GEMMA. Under Bonferroni correction, the number of significant SNPs was heterogeneous across species (Arabidopsis 1160, rice 1536, sorghum 4, barley 125, maize 5). To evaluate cross-species convergence beyond the Bonferroni tail, we selected the top 0.5\% of SNPs per species and mapped nearby genes to orthogroups, using gene-body-only annotation (0 bp) for Arabidopsis and rice and \pm25 kb annotation for maize, sorghum, and barley. This yielded 3,094 Arabidopsis genes, 3,457 rice genes, 3,934 maize genes, 1,153 barley genes, and 940 sorghum genes.

Orthogroup integration identified 20,369 total orthogroups, with 4 orthogroups represented in all five species. Table~\ref{tab:shared_function_candidates_am_roots_colonized_mixed} summarizes these shared orthogroups and reports, for each species, the strongest gene-level signal within each orthogroup as $-\log_{10}(P)$.

\begin{table*}[t]
\caption{Shared orthogroups across all five species from the top 0.5\% of SNPs for AM fungal roots colonized (mixed-window annotation).\label{tab:shared_function_candidates_am_roots_colonized_mixed}}
\scriptsize
\tabcolsep=3pt
\begin{tabular*}{\textwidth}{@{\extracolsep{\fill}}lllllllllll@{}}
\toprule
Orthogroup & Maize gene & $-\log_{10}(P)$ & Sorghum gene & $-\log_{10}(P)$ & Rice gene & $-\log_{10}(P)$ & Barley gene & $-\log_{10}(P)$ & Arabidopsis gene & $-\log_{10}(P)$ \\
\midrule
640662at2759 & Zm00001eb001560 & 4.16 & SORBI\_3001G293000 & 2.45 & LOC\_Os05g01620 & 4.81 & HORVU.MOREX.r3.7HG0667620 & 6.24 & AT1G68850 & 7.51 \\
717421at2759 & Zm00001eb144570 & 2.90 & SORBI\_3004G185000 & 2.85 & LOC\_Os01g08990 & 4.68 & HORVU.MOREX.r3.4HG0338070 & 3.34 & AT2G41970 & 6.56 \\
717421at3193 & Zm00001eb144570 & 2.90 & SORBI\_3004G185000 & 2.85 & LOC\_Os01g08990 & 4.68 & HORVU.MOREX.r3.4HG0338070 & 3.34 & AT2G41970 & 6.56 \\
717421at33090 & Zm00001eb144570 & 2.90 & SORBI\_3004G185000 & 2.85 & LOC\_Os01g08990 & 4.68 & HORVU.MOREX.r3.4HG0338070 & 3.34 & AT2G41970 & 6.56 \\
\botrule
\end{tabular*}
\begin{tablenotes}%
\item Note: One representative gene per species is shown for each orthogroup (the gene with the smallest $P$ within that species/orthogroup).
\end{tablenotes}
\end{table*}

